% !TeX program = xelatex
% ============================================================
% project_description/project_description.tex — «Детальное описание
% проекта» к КТ1/КТ2 курсового проекта ОП «Программная инженерия».
%
% Обязательный деливерабл КТ1 (обновляется на КТ2) по Программе практики:
%   • прикладной вид  — Приложение 6 (планы/этапы, тех.стек+обоснование,
%     >=3 количественных критериев оценки, особые пометки);
%   • исследовательский вид — Приложение 7 (цель, мотивация, объект,
%     предмет, гипотеза, обзор, планы/этапы, >=3 критериев, пометки).
% Структура ветвится по project.kind, язык — по project.lang.
% Заполняйте жёлтые \hseFill{…}; данные о себе и проекте — в настройках.
% ============================================================
\documentclass[12pt,a4paper]{extarticle}
\input{preamble}
\begin{document}
\sloppy

% ── Шапка документа ──
\begin{center}

{\large\bfseries Детальное описание проекта}\\[2pt]
{\normalsize (Контрольные точки 1--2, курсовой проект)}\\[6pt]

\textbf{\projectname}\\[4pt]

\hseAuthorName, группа \hseGroup

\end{center}
\vspace{0.4cm}

((* if project.kind == 'research' *))
% ============================================================
%  ПРИЛОЖЕНИЕ 7 — исследовательский проект
% ============================================================

\section{Основные планы и этапы проекта}
\subsection{Описание проекта}
\textbf{Цель проекта:} \hseFill{одно предложение — какой научный результат достигается}.

\textbf{Мотивация проекта} (почему выбрали именно этот проект): \hseFill{чем тема актуальна и интересна}.

\textbf{Объект исследования:} \hseFill{объект}.

\textbf{Предмет исследования:} \hseFill{предмет}.

\textbf{Краткое описание гипотезы:} \hseFill{проверяемая гипотеза}.

\textbf{Обзор предметной области:} \hseFill{краткий обзор области и существующих подходов}.

\subsection{Планы и этапы выполнения проекта}
\begin{longtable}{|>{\raggedright\arraybackslash}p{3cm}|>{\raggedright\arraybackslash}p{5cm}|>{\raggedright\arraybackslash}p{3.5cm}|>{\raggedright\arraybackslash}p{2cm}|}
\hline
\textbf{Этап проекта} & \textbf{Описание работ} & \textbf{Ожидаемые результаты} & \textbf{Сроки} \\ \hline
\endhead
\hseFill{этап 1} & \hseFill{что делается} & \hseFill{результат} & \hseFill{срок} \\ \hline
\hseFill{этап 2} & \hseFill{что делается} & \hseFill{результат} & \hseFill{срок} \\ \hline
\end{longtable}

\section{Критерии оценивания проекта}
Сформируйте не менее трёх количественных критериев, которыми будет оцениваться проект (например: новизна исследования, качество обзора литературы, корректность теоретических решений, качество программной реализации, глубина анализа результатов экспериментов).
\begin{longtable}{|>{\raggedright\arraybackslash}p{5cm}|>{\raggedright\arraybackslash}p{11cm}|}
\hline
\textbf{Критерий} & \textbf{Описание} \\ \hline
\endhead
\hseFill{критерий 1} & \hseFill{как измеряется} \\ \hline
\hseFill{критерий 2} & \hseFill{как измеряется} \\ \hline
\hseFill{критерий 3} & \hseFill{как измеряется} \\ \hline
\end{longtable}

\section{Особые пометки}
\hseFill{любые примечания, риски, изменения в планах и другие важные детали на данном этапе}.

((* else *))
% ============================================================
%  ПРИЛОЖЕНИЕ 6 — прикладной (программный) проект
% ============================================================

\section{Основные планы и этапы проекта}
\subsection{Краткое описание проекта}
\textbf{Название проекта:} \projectname.

\textbf{Цель проекта:} \hseFill{одно предложение — что создаётся}.

\textbf{Краткое описание задач:} \hseFill{ключевые задачи проекта}.

\subsection{Планы и этапы выполнения проекта}
\begin{longtable}{|>{\raggedright\arraybackslash}p{3cm}|>{\raggedright\arraybackslash}p{5cm}|>{\raggedright\arraybackslash}p{3.5cm}|>{\raggedright\arraybackslash}p{2cm}|}
\hline
\textbf{Этап проекта} & \textbf{Описание работ} & \textbf{Ожидаемые результаты} & \textbf{Сроки} \\ \hline
\endhead
\hseFill{этап 1} & \hseFill{что делается} & \hseFill{результат} & \hseFill{срок} \\ \hline
\hseFill{этап 2} & \hseFill{что делается} & \hseFill{результат} & \hseFill{срок} \\ \hline
\end{longtable}

\section{Используемый технологический стек и его обоснование}
\subsection{Перечень используемых технологий}
\begin{longtable}{|>{\raggedright\arraybackslash}p{4cm}|>{\raggedright\arraybackslash}p{5cm}|>{\raggedright\arraybackslash}p{7cm}|}
\hline
\textbf{Технология / инструмент} & \textbf{Описание} & \textbf{Причины выбора} \\ \hline
\endhead
\hseFill{технология 1} & \hseFill{что это} & \hseFill{почему выбрана} \\ \hline
\hseFill{технология 2} & \hseFill{что это} & \hseFill{почему выбрана} \\ \hline
\end{longtable}

\subsection{Обоснование выбранного технологического стека}
\hseFill{почему выбран именно этот стек, какие преимущества он имеет для проекта и как помогает достичь целей}.

\section{Критерии оценивания проекта}
Сформируйте не менее трёх количественных критериев, которыми будет оцениваться проект (например: \% выполнения функциональных требований, количество реализованных функций, время отклика, покрытие тестами \%, число ошибок на 1000 строк кода, цикломатическая сложность, число поддерживаемых платформ).
\begin{longtable}{|>{\raggedright\arraybackslash}p{5cm}|>{\raggedright\arraybackslash}p{11cm}|}
\hline
\textbf{Критерий} & \textbf{Описание} \\ \hline
\endhead
\hseFill{критерий 1} & \hseFill{как измеряется} \\ \hline
\hseFill{критерий 2} & \hseFill{как измеряется} \\ \hline
\hseFill{критерий 3} & \hseFill{как измеряется} \\ \hline
\end{longtable}

\section{Особые пометки}
\hseFill{любые примечания, риски, изменения в планах и другие важные детали на данном этапе}.

((* endif *))

\end{document}
